\documentclass[xcolor=dvipsnames]{beamer}

\usepackage{amsmath, amssymb, graphicx}
\usepackage[english]{babel}
\usepackage{times}
\usepackage[utf8]{inputenc}
\usepackage[T1]{fontenc}
\usepackage{listings}
\usepackage[norelsize,ruled,vlined]{algorithm2e}
\usepackage{color}
\usepackage{hyperref}
\usepackage{listings}
\usepackage{booktabs}
\usepackage{listings}
\usepackage{tikz}
\usepackage{xcolor}
\usepackage{caption}
\usetikzlibrary{matrix}
\usetikzlibrary{arrows}
\usetikzlibrary{positioning}
\usetikzlibrary{shapes.multipart}

\theoremstyle{definition}
\newtheorem{proposition}{Proposition}
\DeclareCaptionFont{white}{\color{white}}
\DeclareCaptionFormat{listing}{\colorbox[cmyk]{0.43, 0.35, 0.35,0.01}{\parbox{\textwidth}{\hspace{15pt}#1#2#3} } }
\captionsetup[lstlisting]{format=listing, labelfont=white, textfont=white, singlelinecheck=false, margin=0pt, font={bf,footnotesize} }

\newcommand{\fancyh}{\mathcal{H}}
\newcommand{\gangle}[1]{\langle{} #1 \rangle{}}
\newcommand{\myd}{\mathrm{d}}
\newcommand{\NN}{\mathbb{N}}
\newcommand{\RR}{\mathbb{R}}
\newcommand{\ZZ}{\mathbb{Z}}

\definecolor{dkgreen}{rgb}{0,0.6,0}
\definecolor{gray}{rgb}{0.5,0.5,0.5}
\definecolor{mauve}{rgb}{0.58,0,0.82}

\lstset{language=Scala,
  aboveskip=3mm,
  basicstyle={\normalsize\ttfamily\obeyspaces},
  basewidth=0.5em,
  belowskip=3mm,
  gobble=20,
  breakatwhitespace=true,
  breaklines=true,
  columns=fixed,
  commentstyle=\color{dkgreen},
  keywordstyle=\color{blue},
  numbers=none,
  numberstyle=\tiny\color{gray},
  showstringspaces=false,
  stepnumber=1,
  stringstyle=\color{mauve},
  tabsize=2,
}



\mode<presentation>
{\setbeamertemplate{navigation symbols}{}
    \setbeamertemplate{items}[ball]
    % \setbeamertemplate{blocks}[rounded][shadow=true]
    \beamertemplatenavigationsymbolsempty
    \usecolortheme[named=Sepia]{structure}
    \usetheme{Warsaw}
    \useoutertheme{infolines}
    \setbeamercovered{transparent}
}

\definecolor{mygreen}{rgb}{0, 178, 115}

\tikzstyle{block} = [rectangle, draw, font=\tiny, text centered, rounded corners, minimum height=1em, node distance=1.5cm, minimum width=2em]
\tikzstyle{line} = [draw, -latex']
\tikzstyle{value} = [label, black, font=\tiny, thick, node distance=0.4cm]

\title[Scalafmt]{Scalafmt \\ {\small Pretty-printing Scala}}
\author{Ólafur Páll Geirsson \\
{\tiny olafur.geirsson@epfl.ch}}
% \department{School of Computer and Communication Sciences}
% \project{Semester Project}
\institute[EPFL]{École Polytechnique Fédérale de Lausanne \\
    School of Computer and Communication Sciences\\
\logoepfl}
\date{\today}


% Delete this, if you do not want the table of contents to pop up at
% the beginning of each subsection:
\AtBeginSection[]
{\begin{frame}<beamer>{Overview}
        \tableofcontents[
            sections={1-6},
            currentsection,
            currentsubsection,
            hideothersubsections,
            sectionstyle=show/shaded,
        subsectionstyle=show/shaded/hide]
    \end{frame}
}

\begin{document}

\newcommand{\logoepfl}[0]{\begin{center}
    \includegraphics[width=2cm]{logo_epfl_coul.eps}
    \vspace{-0.5cm}
  \end{center}
}


%%%%%%%%%%%%%%%%%%%%%%%%%%%%%%%%%%%%%%%%%%%%%%%%%%%%%%%%%%%%%%%
% 0. Titlepage
%%%%%%%%%%%%%%%%%%%%%%%%%%%%%%%%%%%%%%%%%%%%%%%%%%%%%%%%%%%%%%%
\begin{frame}
    \titlepage{}
\end{frame}

%%%%%%%%%%%%%%%%%%%%%%%%%%%%%%%%%%%%%%%%%%%%%%%%%%%%%%%%%%%%%%%
% 1. Motivation
%%%%%%%%%%%%%%%%%%%%%%%%%%%%%%%%%%%%%%%%%%%%%%%%%%%%%%%%%%%%%%%
\section{Motivation} % (fold)
\label{sec:Motivation}
\begin{frame}[fragile]
    \frametitle{Motivation}
    \begin{block}{}
        \begin{quote}
            {\large ``any style guide written in English is either so brief that it’s ambiguous, or so long that no one reads it.''}
            \vskip2mm
            \hspace*\fill{\small--- Bob Nystrom, Dart team}
        \end{quote}
    \end{block}
\end{frame}
\begin{frame}[fragile]
    \frametitle{Motivation}
    \begin{block}{The dream}
        \includegraphics[width=\linewidth]{img/klang.png}
    \end{block}
\end{frame}
\begin{frame}[fragile]
    \frametitle{Motivation}
    \begin{block}{}
        \begin{itemize}
            \item Code review.
            \item Consistent style.
            \item Refactoring.
        \end{itemize}
    \end{block}
\end{frame}

%%%%%%%%%%%%%%%%%%%%%%%%%%%%%%%%%%%%%%%%%%%%%%%%%%%%%%%%%%%%%%%
% 2. Scalafmt
%%%%%%%%%%%%%%%%%%%%%%%%%%%%%%%%%%%%%%%%%%%%%%%%%%%%%%%%%%%%%%%
\section{Scalafmt} % (fold)
\label{sec:Scalafmt}

\begin{frame}[fragile]
    \frametitle{Interface}
    \begin{block}{}
        \lstinputlisting{code/style.scala}
        \lstinputlisting{code/interface.scala}
    \end{block}
\end{frame}

\begin{frame}[fragile]
    \frametitle{Multiple arguments}
    \begin{block}{}
        \lstinputlisting{code/example-apply.scala}
    \end{block}
\end{frame}

\begin{frame}[fragile]
    \frametitle{Inline if assignment}
    \begin{block}{}
        \lstinputlisting{code/example-if.scala}
    \end{block}
\end{frame}

\begin{frame}[fragile]
    \frametitle{Pattern matching}
    \begin{block}{}
        \lstinputlisting{code/example-case.scala}
    \end{block}
\end{frame}

%%%%%%%%%%%%%%%%%%%%%%%%%%%%%%%%%%%%%%%%%%%%%%%%%%%%%%%%%%%%%%%
% 3. Conclusion
%%%%%%%%%%%%%%%%%%%%%%%%%%%%%%%%%%%%%%%%%%%%%%%%%%%%%%%%%%%%%%%
\section{Conclusion} % (fold)
\label{sec:Conclusion}

\begin{frame}[fragile]
    \frametitle{Conclusion}
    \begin{block}{Scalafmt}
        \begin{itemize}
            \item ???
        \end{itemize}
    \end{block}
\end{frame}

\begin{frame}[fragile]
    \frametitle{The End}
    \begin{center}
        \Huge
        Thank you!
    \end{center}
\end{frame}

% subsection Evaluation (end)


\end{document}
